\chapter{Jordan 표준형}

\chapterintro{이 장에서는 대각화가 되지 않는 선형연산자까지 구조적으로 분류하는 Jordan 표준형(Jordan canonical form)을 다룹니다. 앞 장에서 준비한 primary 분해와 cyclic 분해를 결합해 존재정리를 완성하고, 핵 차원 증가량을 이용해 유일성을 확인하겠습니다. 마지막에는 $A^n$, $e^{tA}$ 계산으로 연결하여 실제 계산 도구로 정리하겠습니다.}

\chaptergoals{Jordan 블록, 일반화 고유벡터, Jordan 사슬의 관계를 정확히 설명할 수 있다.}{대수적으로 닫힌 체에서 Jordan 표준형의 존재성과 유일성을 완전한 논리로 설명할 수 있다.}{Jordan 표준형을 이용하여 거듭제곱과 지수행렬 계산을 체계적으로 수행할 수 있다.}

\sectiontemplate{Jordan 블록과 Jordan 사슬}{Jordan 이론의 출발점은 한 개의 블록을 정확히 이해하는 것입니다. 이 절에서는 블록, 일반화 고유벡터, 사슬을 정의하고, 사슬이 왜 블록 행렬을 만드는지 확인하겠습니다.}{\item Jordan 블록과 Jordan 사슬의 정의를 말할 수 있습니다.\item 사슬 기저에서 행렬이 Jordan 블록이 되는 이유를 설명할 수 있습니다.\item Jordan 블록의 핵 차원 증가량을 계산할 수 있습니다.}{\item 고유값과 고유벡터\item 불변부분공간\item 최소다항식과 특성다항식}

\begin{definition}
$\lambda\in F$, $k\ge 1$에 대하여
$$
J_k(\lambda)=
\begin{bmatrix}
\lambda&1&0&\cdots&0\\
0&\lambda&1&\ddots&\vdots\\
\vdots&\ddots&\ddots&\ddots&0\\
\vdots&&\ddots&\lambda&1\\
0&\cdots&\cdots&0&\lambda
\end{bmatrix}
\in M_k(F)
$$
를 크기 $k$의 Jordan 블록(Jordan block)이라 한다.
\end{definition}

\begin{definition}
$T\in\operatorname{End}(V)$, $\lambda\in F$에 대하여, $v\neq 0$이고 어떤 $r\ge 1$에 대해 $(T-\lambda I)^r v=0$이면 $v$를 $\lambda$에 대한 일반화 고유벡터(generalized eigenvector)라 한다. 또한 벡터열 $(v_1,\dots,v_k)$가
$$
(T-\lambda I)v_1=0,\qquad (T-\lambda I)v_{j}=v_{j-1}\ \ (2\le j\le k)
$$
를 만족하면 길이 $k$의 Jordan 사슬(Jordan chain)이라 한다.
\end{definition}

\begin{theorem}[Jordan 사슬의 행렬표현]
$T\in\operatorname{End}(V)$, $\lambda\in F$라 하자. $(v_1,\dots,v_k)$가 $\lambda$에 대한 Jordan 사슬이면 $W=\operatorname{span}\{v_1,\dots,v_k\}$는 $T$-불변부분공간이며, 기저 $\mathcal{B}=(v_1,\dots,v_k)$에 대한 $T|_W$의 행렬은 $J_k(\lambda)$이다.
\end{theorem}

\paragraph{증명 전략.}
사슬의 점화식에서 각 $Tv_j$를 기저 $\mathcal{B}$ 좌표로 직접 적습니다. 그러면 열벡터가 Jordan 블록의 열과 정확히 일치함을 확인할 수 있습니다.

\begin{proof}
사슬 조건에서
$$
Tv_1=\lambda v_1,\qquad Tv_j=\lambda v_j+v_{j-1}\ \ (2\le j\le k)
$$
가 성립한다. 따라서 각 $Tv_j$는 $v_1,\dots,v_k$의 선형결합이므로 $W$는 $T$-불변이다.

이제 기저 $\mathcal{B}$에서의 좌표를 본다. $Tv_1$의 좌표열은 $(\lambda,0,\dots,0)^T$이고, $j\ge2$에 대해 $Tv_j$의 좌표열은 $(0,\dots,0,1,\lambda,0,\dots,0)^T$이며 $1$은 $(j-1)$번째, $\lambda$는 $j$번째 위치에 놓인다. 그러므로 $[T|_W]_{\mathcal{B}}$는 대각에 $\lambda$, 초대각에 $1$, 나머지에 $0$을 갖는다. 즉
$$
[T|_W]_{\mathcal{B}}=J_k(\lambda).
$$
\end{proof}

\paragraph{의미 해석.}
Jordan 사슬 하나는 Jordan 블록 하나와 정확히 대응합니다. 따라서 Jordan 표준형 문제는 결국 ``어떤 길이의 사슬들이 몇 개 필요한가''를 결정하는 문제로 바뀝니다.

\begin{theorem}[Jordan 블록의 핵 차원과 최소다항식]
$J=J_k(\lambda)$, $N=J-\lambda I$라 하자. 임의의 $m\ge 1$에 대하여
$$
\dim\ker (J-\lambda I)^m=\min\{m,k\}
$$
가 성립한다. 또한
$$
\chi_J(t)=(t-\lambda)^k,\qquad m_J(t)=(t-\lambda)^k
$$
이다.
\end{theorem}

\paragraph{증명 전략.}
$N$은 초대각 이동 연산자이므로 표준기저에 대한 작용이 명시적입니다. $N^m$의 핵을 직접 계산하고, $N^{k-1}\neq0$, $N^k=0$으로 최소다항식 차수를 결정합니다.

\begin{proof}
표준기저 $e_1,\dots,e_k$에 대해
$$
Ne_1=0,\qquad Ne_i=e_{i-1}\ \ (2\le i\le k)
$$
이다. 귀납적으로
$$
N^m e_i=
\begin{cases}
e_{i-m},& i>m,\\
0,& i\le m
\end{cases}
$$
를 얻는다. 따라서 $m<k$이면 $\ker N^m=\operatorname{span}\{e_1,\dots,e_m\}$이고, $m\ge k$이면 $\ker N^m=F^k$이다. 즉
$$
\dim\ker (J-\lambda I)^m=\dim\ker N^m=\min\{m,k\}.
$$

다음으로 $J$는 상삼각행렬이며 대각 원소가 모두 $\lambda$이므로 특성다항식은
$$
\chi_J(t)=(t-\lambda)^k.
$$
또한 $N^k=0$이므로 $(J-\lambda I)^k=0$, 따라서 최소다항식은 $(t-\lambda)^k$를 나눈다. 한편
$$
N^{k-1}e_k=e_1\neq0
$$
이므로 $(J-\lambda I)^{k-1}\neq0$이다. 따라서 최소다항식의 지수는 정확히 $k$이고
$$
m_J(t)=(t-\lambda)^k
$$
이다.
\end{proof}

\paragraph{의미 해석.}
핵 차원 증가량은 사슬 길이를 읽는 핵심 데이터입니다. 이후 유일성 정리에서 $\dim\ker(T-\lambda I)^m$이 블록 크기를 복원하는 불변량으로 작동합니다.

\begin{example}
$J_4(3)$에 대해 $N=J_4(3)-3I$이면
$$
\dim\ker N=1,\quad \dim\ker N^2=2,\quad \dim\ker N^3=3,\quad \dim\ker N^4=4.
$$
또한 $m_{J_4(3)}(t)=(t-3)^4$이다.
\end{example}

\subsection*{주의}
\begin{warning}
Jordan 사슬을 $(v_k,\dots,v_1)$ 순서로 쓰면 초대각의 $1$ 위치가 달라집니다. 기저 순서를 고정하고 계산해야 실수가 줄어듭니다.
\end{warning}
\begin{warning}
$(T-\lambda I)^r v=0$이라고 해서 곧바로 사슬이 주어지지는 않습니다. 사슬은 각 단계가 이전 벡터를 정확히 만들어내는 추가 조건을 포함합니다.
\end{warning}

\subsection*{자가진단퀴즈}
\begin{enumerate}[label=\arabic*.]
\item $J_5(-2)$에 대해 $\dim\ker(J_5(-2)+2I)^m$을 $m=1,2,3,4,5$에 대해 구하십시오.
\item 길이 $3$의 Jordan 사슬 $(v_1,v_2,v_3)$에서 $Tv_2$, $Tv_3$를 $\{v_1,v_2,v_3\}$ 좌표로 쓰고, 해당 열벡터를 적으십시오.
\item $J_k(\lambda)$의 최소다항식 지수가 왜 $k$보다 작아질 수 없는지 한 줄 논증으로 설명하십시오.
\end{enumerate}

\sectiontemplate{Nilpotent 연산자의 Jordan 분해}{일반 Jordan 정리를 증명하기 전에, 먼저 고유값이 $0$뿐인 nilpotent 경우를 완성하면 구조가 크게 단순해집니다. 이 절에서 존재성과 유일성을 모두 정리하겠습니다.}{\item nilpotent 연산자의 Jordan 분해 존재정리를 설명할 수 있습니다.\item 핵 차원 증가량으로 블록 크기를 복원할 수 있습니다.\item 분할(partition) 데이터와 Jordan 블록의 대응을 계산할 수 있습니다.}{\item cyclic subspace\item direct sum 분해\item nilpotent 지수}

\begin{definition}
$N\in\operatorname{End}(V)$에 대해 어떤 $s\ge1$가 존재하여 $N^s=0$이면 $N$을 nilpotent 연산자라 한다. 최소의 그런 $s$를 $N$의 nilpotent 지수(nilpotency index)라 한다.
\end{definition}

\begin{theorem}[Nilpotent Jordan 분해 존재정리]
$V$가 유한차원이고 $N\in\operatorname{End}(V)$가 nilpotent이면, 어떤 기저에서 $N$의 행렬은
$$
J_{k_1}(0)\oplus\cdots\oplus J_{k_s}(0)
$$
꼴이다.
\end{theorem}

\paragraph{증명 전략.}
앞 장의 cyclic 분해정리를 사용해 $V$를 $N$-cyclic 부분공간의 직합으로 나눕니다. 각 cyclic 부분공간에서는 적절한 사슬 기저를 택하면 행렬이 단일 Jordan 블록이 됨을 보입니다.

\begin{proof}
앞 장의 cyclic decomposition 정리에 의해
$$
V=Z(v_1;N)\oplus\cdots\oplus Z(v_s;N)
$$
가 되도록 $v_1,\dots,v_s$를 잡을 수 있다. 각 $i$에 대해
$$
k_i=\min\{k\ge1: N^k v_i=0\}
$$
로 둔다. $N$이 nilpotent이므로 $k_i$는 유한하다.

$W_i=Z(v_i;N)$에서 벡터열
$$
\mathcal{B}_i=(N^{k_i-1}v_i,\ N^{k_i-2}v_i,\ \dots,\ Nv_i,\ v_i)
$$
를 보자. 이 벡터열이 $W_i$의 기저임을 보이면 충분하다.

우선 생성성은 $W_i=\operatorname{span}\{v_i,Nv_i,\dots\}$ 정의에서 자명하다. 선형독립성을 보이자. 만약
$$
a_0 v_i+a_1 Nv_i+\cdots+a_{k_i-1}N^{k_i-1}v_i=0
$$
이면 다항식 $p(t)=a_0+a_1 t+\cdots+a_{k_i-1}t^{k_i-1}$에 대해 $p(N)v_i=0$이다. 그런데 $k_i$의 최소성 때문에 $v_i$를 소거하는 최소 차수 단항식은 $t^{k_i}$이므로 $\deg p<k_i$인 $p$는 영다항식이어야 한다. 따라서 $a_0=\cdots=a_{k_i-1}=0$이고 선형독립이다.

이제 $N$의 작용을 쓰면
$$
N(N^{k_i-1}v_i)=0,\qquad N(N^{k_i-j}v_i)=N^{k_i-j+1}v_i\ (2\le j\le k_i),
$$
즉 기저 $\mathcal{B}_i$에서 행렬은 $J_{k_i}(0)$이다. 마지막으로 $V$가 $W_i$들의 직합이므로 $\mathcal{B}_1,\dots,\mathcal{B}_s$를 이어 붙인 기저에서
$$
[N]=J_{k_1}(0)\oplus\cdots\oplus J_{k_s}(0)
$$
를 얻는다.
\end{proof}

\paragraph{의미 해석.}
nilpotent 경우는 Jordan 이론의 핵심 원형입니다. 일반 연산자도 각 고유값 성분에서 ``$\lambda I+$nilpotent'' 꼴로 환원되므로, 이 정리가 전체 존재정리의 엔진이 됩니다.

\begin{theorem}[Nilpotent Jordan 분해 유일성]
$N\sim J_{k_1}(0)\oplus\cdots\oplus J_{k_s}(0)$라 하자. $j\ge1$에 대해
$$
b_j=\dim\ker N^j-\dim\ker N^{j-1}\quad(\ker N^0=\{0\})
$$
로 두면
$$
b_j=\#\{i: k_i\ge j\}
$$
가 성립한다. 특히 블록 크기의 다중집합 $\{k_1,\dots,k_s\}$는 $N$으로부터 유일하게 결정된다.
\end{theorem}

\paragraph{증명 전략.}
블록별로 $\dim\ker J_{k_i}(0)^j=\min\{j,k_i\}$를 합산합니다. 차분을 취하면 ``길이가 $j$ 이상인 블록 개수''가 나타나고, 다시 차분하면 길이가 정확히 $j$인 블록 개수를 복원할 수 있습니다.

\begin{proof}
유사한 행렬은 $\dim\ker N^j$를 보존하므로 Jordan 블록 합으로 계산해도 된다. 각 블록에 대해
$$
\dim\ker J_{k_i}(0)^j=\min\{j,k_i\}
$$
이므로
$$
\dim\ker N^j=\sum_{i=1}^s \min\{j,k_i\}.
$$
따라서
$$
b_j=\sum_{i=1}^s\left(\min\{j,k_i\}-\min\{j-1,k_i\}\right).
$$
괄호 안은 $k_i\ge j$일 때 $1$, 아니면 $0$이므로
$$
b_j=\#\{i:k_i\ge j\}.
$$

이제 $c_j=\#\{i:k_i=j\}$라 두면
$$
c_j=b_j-b_{j+1}
$$
이다. $b_j$들은 $\dim\ker N^j$로부터 전부 결정되므로 $c_j$도 전부 결정된다. 즉 각 길이의 블록 개수가 유일하고, 결과적으로 블록 크기 다중집합이 유일하다.
\end{proof}

\paragraph{의미 해석.}
유일성의 본질은 ``핵 차원 증가량 = 사슬 개수''라는 대응입니다. Jordan 형식을 기억할 때 블록 자체를 외우기보다 $\dim\ker N^j$ 표를 먼저 만드는 습관이 계산에 강합니다.

\begin{example}
$\dim V=7$인 nilpotent $N$에 대해
$$
\dim\ker N=3,\quad \dim\ker N^2=5,\quad \dim\ker N^3=6,\quad \dim\ker N^4=7
$$
라 하자. 그러면
$$
b_1=3,\ b_2=2,\ b_3=1,\ b_4=1
$$
이므로
$$
c_1=b_1-b_2=1,\quad c_2=b_2-b_3=1,\quad c_3=b_3-b_4=0,\quad c_4=b_4-b_5=1
$$
($b_5=0$)가 되어 블록 크기는 $4,2,1$이다.
\end{example}

\subsection*{주의}
\begin{warning}
$\dim\ker N^j$ 값만 보고 바로 블록 크기를 추정하면 자주 틀립니다. 반드시 차분 $b_j=\dim\ker N^j-\dim\ker N^{j-1}$을 먼저 계산해야 합니다.
\end{warning}
\begin{warning}
nilpotent 지수와 블록 개수는 다른 정보입니다. 지수는 ``최대 블록 크기''만 알려 주며, 전체 분할은 핵 차원 데이터가 필요합니다.
\end{warning}

\subsection*{자가진단퀴즈}
\begin{enumerate}[label=\arabic*.]
\item $\dim\ker N=2,\ \dim\ker N^2=4,\ \dim\ker N^3=5,\ \dim\ker N^4=5$일 때 Jordan 블록 크기를 구하십시오.
\item 왜 $b_j$가 항상 감소수열($b_{j+1}\le b_j$)인지 설명하십시오.
\item nilpotent 연산자에서 최소다항식의 차수가 최대 Jordan 블록 크기와 일치함을 증명하십시오.
\end{enumerate}

\sectiontemplate{일반 연산자의 Jordan 표준형: 존재와 유일성}{이제 일반 연산자로 확장합니다. 핵심은 고유값별 일반화 고유공간으로 쪼갠 뒤, 각 성분에서 nilpotent 정리를 적용하는 것입니다.}{\item 일반화 고유공간의 직합분해를 사용할 수 있습니다.\item Jordan 표준형 존재정리를 완전한 구조로 설명할 수 있습니다.\item 유일성의 불변량을 고유값별로 계산할 수 있습니다.}{\item primary decomposition\item 고유값 분해\item nilpotent Jordan 정리}

\begin{definition}
$T\in\operatorname{End}(V)$의 최소다항식이
$$
m_T(t)=\prod_{i=1}^r (t-\lambda_i)^{e_i}
$$
($\lambda_i$ 서로 다름)로 분해된다고 하자. 각 $i$에 대해
$$
G_{\lambda_i}(T)=\ker (T-\lambda_i I)^{e_i}
$$
를 $\lambda_i$에 대한 일반화 고유공간(generalized eigenspace)이라 한다.
\end{definition}

\begin{theorem}[일반화 고유공간의 직합분해]
위 표기에서
$$
V=G_{\lambda_1}(T)\oplus\cdots\oplus G_{\lambda_r}(T)
$$
가 성립한다.
\end{theorem}

\paragraph{증명 전략.}
primary decomposition 정리를 최소다항식의 서로소 인수 $(t-\lambda_i)^{e_i}$에 직접 적용합니다. 정의를 대입하면 각 항이 일반화 고유공간이 됩니다.

\begin{proof}
다항식
$$
p_i(t)=(t-\lambda_i)^{e_i}\quad(1\le i\le r)
$$
를 두면 $p_i$들은 서로소이며 $m_T(t)=p_1(t)\cdots p_r(t)$이다. primary decomposition 정리에 의해
$$
V=\ker p_1(T)\oplus\cdots\oplus\ker p_r(T).
$$
그런데 $\ker p_i(T)=\ker (T-\lambda_i I)^{e_i}=G_{\lambda_i}(T)$이므로 결론이 따른다.
\end{proof}

\paragraph{의미 해석.}
Jordan 문제를 고유값별 독립 문제로 분해할 수 있게 됩니다. 이후 계산은 각 $G_{\lambda_i}(T)$에서 별도로 수행하고 마지막에 직합으로 합치면 됩니다.

\begin{theorem}[Jordan 표준형 존재정리]
$F$가 대수적으로 닫힌 체이고 $V$가 유한차원일 때, 임의의 $T\in\operatorname{End}(V)$에 대해 어떤 기저에서
$$
[T]=\bigoplus_{i=1}^r\ \bigoplus_{\ell=1}^{s_i} J_{k_{i,\ell}}(\lambda_i)
$$
가 된다.
\end{theorem}

\paragraph{증명 전략.}
직합분해 $V=\oplus_i G_{\lambda_i}(T)$를 먼저 만든 뒤, 각 성분에서 $N_i=(T-\lambda_i I)|_{G_{\lambda_i}(T)}$를 nilpotent로 확인합니다. 그러면 이전 절의 nilpotent 존재정리를 성분별로 적용할 수 있습니다.

\begin{proof}
앞 정리에 의해
$$
V=\bigoplus_{i=1}^r G_{\lambda_i}(T).
$$
각 $i$에 대해 $G_i=G_{\lambda_i}(T)$, $N_i=(T-\lambda_i I)|_{G_i}$라 두자. $x\in G_i$이면 정의상 $(T-\lambda_i I)^{e_i}x=0$이므로 $N_i^{e_i}=0$, 따라서 $N_i$는 nilpotent이다.

이제 nilpotent Jordan 분해 존재정리를 $N_i$에 적용하면 $G_i$의 어떤 기저에서
$$
[N_i]=\bigoplus_{\ell=1}^{s_i} J_{k_{i,\ell}}(0).
$$
그러면 같은 기저에서
$$
[T|_{G_i}]=\lambda_i I+[N_i]
=\bigoplus_{\ell=1}^{s_i} J_{k_{i,\ell}}(\lambda_i).
$$
마지막으로 $G_1,\dots,G_r$의 기저를 이어 붙이면 $V$의 기저가 되고, 그 기저에서 $[T]$는 위 블록 직합 형태가 된다.
\end{proof}

\paragraph{의미 해석.}
일반 Jordan 존재정리는 ``고유값별 분해 + nilpotent 표준화''의 합성입니다. 즉 새로운 아이디어는 많지 않고, 앞 장과 앞 절의 결과를 정확히 조립하는 것이 핵심입니다.

\begin{theorem}[Jordan 표준형 유일성(블록 순서 제외)]
$A,B\in M_n(F)$가 서로 유사하고 둘 다 Jordan 형태라 하자. 그러면 각 고유값 $\lambda$에 대해 $A$와 $B$의 $\lambda$-Jordan 블록 크기 다중집합은 일치한다.
\end{theorem}

\paragraph{증명 전략.}
고유값 $\lambda$를 고정하고 $d_{\lambda,j}=\dim\ker(A-\lambda I)^j$를 불변량으로 사용합니다. $\mu\neq\lambda$인 블록은 핵에 기여하지 않음을 확인한 뒤, nilpotent 유일성 정리와 같은 차분 논법으로 블록 크기를 복원합니다.

\begin{proof}
고유값 $\lambda$를 고정하고
$$
d_{\lambda,j}(A)=\dim\ker(A-\lambda I)^j,\qquad
d_{\lambda,j}(B)=\dim\ker(B-\lambda I)^j
$$
로 두자. 유사성 $B=P^{-1}AP$에서
$$
(B-\lambda I)^j=P^{-1}(A-\lambda I)^jP
$$
이므로
$$
d_{\lambda,j}(A)=d_{\lambda,j}(B)\quad(\forall j\ge1).
$$

이제 Jordan 형태에서 $A$의 블록 기여를 본다. $\mu\neq\lambda$인 블록 $J_k(\mu)$에 대해
$$
J_k(\mu)-\lambda I=(\mu-\lambda)I+N_k
$$
($N_k$ nilpotent). $\mu-\lambda\neq0$이므로
$$
\big((\mu-\lambda)I+N_k\big)^{-1}
=(\mu-\lambda)^{-1}\sum_{r=0}^{k-1}\left(-(\mu-\lambda)^{-1}N_k\right)^r
$$
가 되어 가역이다. 따라서 그 거듭제곱도 가역이므로 핵 기여가 0이다.

반면 $\lambda$-블록 $J_{k}(\lambda)$는
$$
\dim\ker(J_k(\lambda)-\lambda I)^j=\min\{j,k\}
$$
를 기여한다. 따라서
$$
d_{\lambda,j}(A)=\sum_{\text{$\lambda$-블록 }k}\min\{j,k\}.
$$
같은 식이 $B$에도 성립하고 $d_{\lambda,j}(A)=d_{\lambda,j}(B)$이므로, 차분
$$
b_{\lambda,j}=d_{\lambda,j}-d_{\lambda,j-1}
$$
은 ``크기가 $j$ 이상인 $\lambda$-블록 개수''를 동일하게 준다. 다시
$$
c_{\lambda,j}=b_{\lambda,j}-b_{\lambda,j+1}
$$
를 취하면 ``크기가 정확히 $j$인 $\lambda$-블록 개수''가 복원된다. 따라서 $A$와 $B$의 $\lambda$-블록 크기 다중집합은 일치한다. $\lambda$가 임의였으므로 정리가 성립한다.
\end{proof}

\paragraph{의미 해석.}
Jordan 형식의 유일성은 행렬 자체의 유일성이 아니라 블록 분할의 유일성입니다. 계산 실무에서는 ``고유값별 핵 차원 표''를 만들면 유일성 확인이 자동으로 해결됩니다.

\begin{example}
어떤 $A\in M_7(F)$가
$$
\chi_A(t)=(t-2)^5(t+1)^2
$$
를 만족하고
$$
\dim\ker(A-2I)=2,\ \dim\ker(A-2I)^2=4,\ \dim\ker(A-2I)^3=5,
$$
$$
\dim\ker(A+I)=1,\ \dim\ker(A+I)^2=2
$$
라 하자. 그러면 $\lambda=2$에 대해 증가량은 $2,2,1$이므로 블록 크기는 $3,2$, $\lambda=-1$에 대해 증가량은 $1,1$이므로 블록 크기는 $2$이다. 따라서 Jordan 형식은
$$
J_3(2)\oplus J_2(2)\oplus J_2(-1)
$$
이다.
\end{example}

\subsection*{주의}
\begin{warning}
Jordan 표준형 존재정리는 체가 대수적으로 닫혀 있어야 합니다. 실수체에서 복소 고유값이 생기면 실수 Jordan 형식만으로는 완전 분류가 되지 않습니다.
\end{warning}
\begin{warning}
유일성에서 ``블록 배열 순서''까지 고정된다고 생각하면 잘못입니다. 순서는 바뀔 수 있고, 블록 크기 다중집합만 불변입니다.
\end{warning}

\subsection*{자가진단퀴즈}
\begin{enumerate}[label=\arabic*.]
\item $m_T(t)=(t-1)^3(t+2)^2$일 때 일반화 고유공간이 어떻게 분해되는지 쓰십시오.
\item $\mu\neq\lambda$에서 $J_k(\mu)-\lambda I$가 왜 가역인지 유한급수 역행렬을 사용해 설명하십시오.
\item $\dim\ker(T-\lambda I)^j$ 데이터만으로 블록 크기를 복원하는 절차를 세 단계로 정리하십시오.
\end{enumerate}

\sectiontemplate{계산 절차와 행렬함수 응용}{Jordan 표준형은 분류 도구이면서 동시에 계산 도구입니다. 이 절에서는 대각화 판정, 거듭제곱, 지수행렬 공식을 정리해 실제 계산에 바로 연결하겠습니다.}{\item 대각화 가능성을 Jordan 블록 크기로 판정할 수 있습니다.\item $J_k(\lambda)^n$과 $e^{tJ_k(\lambda)}$를 닫힌식으로 계산할 수 있습니다.\item $A=PJP^{-1}$를 이용한 행렬함수 계산 절차를 적용할 수 있습니다.}{\item Jordan 존재/유일성\item 이항정리\item 멱급수와 지수행렬}

\begin{definition}
$A\in M_n(F)$가 가역행렬 $P$와 Jordan 행렬 $J$에 대해
$$
A=PJP^{-1}
$$
로 표현될 때, $J$를 $A$의 Jordan 표준형이라 하고 $P$의 열벡터 기저를 Jordan 기저(Jordan basis)라 한다.
\end{definition}

\begin{theorem}[대각화 가능성 판정]
$F$가 대수적으로 닫혀 있고 $T\in\operatorname{End}(V)$라 하자. 다음은 서로 동치이다.
\begin{enumerate}[label=(\arabic*)]
\item $T$는 대각화 가능하다.
\item $T$의 모든 Jordan 블록 크기가 $1$이다.
\item $m_T(t)$는 중근이 없다.
\end{enumerate}
\end{theorem}

\paragraph{증명 전략.}
Jordan 블록 관점에서 (1)$\Leftrightarrow$(2)를 직접 확인하고, 최소다항식이 각 고유값 블록의 최대 길이를 지수로 갖는다는 사실로 (2)$\Leftrightarrow$(3)을 연결합니다.

\begin{proof}
(1)$\Rightarrow$(2): $T$가 대각화 가능하면 어떤 기저에서 대각행렬이 된다. 대각행렬의 Jordan 블록은 모두 크기 $1$이다.

(2)$\Rightarrow$(1): Jordan 블록이 전부 크기 $1$이면 Jordan 행렬 자체가 대각행렬이므로 $T$는 대각화 가능하다.

(2)$\Rightarrow$(3): Jordan 형식에서 최소다항식은 각 고유값 $\lambda$에 대해 ``$\lambda$-블록의 최대 크기''를 지수로 갖는
$$
m_T(t)=\prod_{\lambda}(t-\lambda)^{s_\lambda}
$$
형태이다. 모든 블록 크기가 $1$이면 모든 $s_\lambda=1$이므로 중근이 없다.

(3)$\Rightarrow$(2): 반대로 어떤 $\lambda$에 대해 크기 $\ge2$ 블록이 있으면 $(t-\lambda)^2$가 최소다항식을 나누어 중근이 생긴다. 이는 (3)에 모순이다. 따라서 모든 블록 크기는 $1$이다.

세 조건이 동치이다.
\end{proof}

\paragraph{의미 해석.}
대각화 실패의 원인은 ``고유값 부족''이 아니라 ``사슬 길이 2 이상 블록의 존재''입니다. 최소다항식의 중근 여부는 이 현상을 한 번에 판정하는 압축된 신호입니다.

\begin{theorem}[Jordan 블록의 거듭제곱과 지수행렬]
$J=J_k(\lambda)=\lambda I+N$ ($N^k=0$)라 하자.
\begin{enumerate}[label=(\arabic*)]
\item 임의의 정수 $n\ge0$에 대해
$$
J^n=\sum_{r=0}^{\min\{n,k-1\}}\binom{n}{r}\lambda^{\,n-r}N^r.
$$
\item 임의의 스칼라 $t$에 대해
$$
e^{tJ}=e^{\lambda t}\sum_{r=0}^{k-1}\frac{t^r}{r!}N^r.
$$
\end{enumerate}
\end{theorem}

\paragraph{증명 전략.}
$\lambda I$와 $N$이 가환함을 이용해 이항정리를 적용합니다. 지수행렬은 멱급수 정의를 사용하고, $N^k=0$ 덕분에 유한합으로 절단됨을 보입니다.

\begin{proof}
(1) $\lambda I$와 $N$은 가환하므로
$$
J^n=(\lambda I+N)^n=\sum_{r=0}^n \binom{n}{r}\lambda^{\,n-r}N^r.
$$
그런데 $N^r=0$ for $r\ge k$이므로
$$
J^n=\sum_{r=0}^{\min\{n,k-1\}}\binom{n}{r}\lambda^{\,n-r}N^r.
$$

(2) 지수행렬 정의로부터
$$
e^{tJ}=\sum_{m=0}^\infty \frac{t^m}{m!}J^m
=\sum_{m=0}^\infty \frac{t^m}{m!}\sum_{r=0}^{\min\{m,k-1\}}\binom{m}{r}\lambda^{\,m-r}N^r.
$$
$N^r$는 $r=0,\dots,k-1$만 남으므로 합 순서를 바꿔
$$
e^{tJ}
=\sum_{r=0}^{k-1}\left(\sum_{m=r}^\infty \frac{t^m}{m!}\binom{m}{r}\lambda^{\,m-r}\right)N^r.
$$
$m=r+q$로 두면
$$
\sum_{m=r}^\infty \frac{t^m}{m!}\binom{m}{r}\lambda^{\,m-r}
=\frac{t^r}{r!}\sum_{q=0}^\infty \frac{(\lambda t)^q}{q!}
=\frac{t^r}{r!}e^{\lambda t}.
$$
따라서
$$
e^{tJ}=e^{\lambda t}\sum_{r=0}^{k-1}\frac{t^r}{r!}N^r.
$$
\end{proof}

\paragraph{의미 해석.}
Jordan 형식으로 바꾸면 고차 행렬함수 계산이 ``대각지수 $e^{\lambda t}$''와 ``유한 다항식''의 곱으로 분해됩니다. 이 때문에 미분방정식 해석과 장기 거듭제곱 계산이 매우 단순해집니다.

\begin{example}
$$
J_3(1)=
\begin{bmatrix}
1&1&0\\
0&1&1\\
0&0&1
\end{bmatrix}
=I+N,\qquad
N=
\begin{bmatrix}
0&1&0\\
0&0&1\\
0&0&0
\end{bmatrix},\quad N^3=0.
$$
정리에 의해
$$
J_3(1)^n=
\begin{bmatrix}
1&n&\binom{n}{2}\\
0&1&n\\
0&0&1
\end{bmatrix},
\qquad
e^{tJ_3(1)}=e^t
\begin{bmatrix}
1&t&\frac{t^2}{2}\\
0&1&t\\
0&0&1
\end{bmatrix}.
$$
\end{example}

\subsection*{주의}
\begin{warning}
수치해석에서는 Jordan 분해가 매우 불안정할 수 있습니다. 이론적 분류와 수치 알고리즘의 안정성은 별개라는 점을 구분해야 합니다.
\end{warning}
\begin{warning}
$A=PJP^{-1}$ 계산에서 $P$를 잘못 정규화하면 이후 $A^n$, $e^{tA}$가 모두 틀립니다. 반드시 $AP=PJ$를 역대입으로 검산해야 합니다.
\end{warning}

\subsection*{자가진단퀴즈}
\begin{enumerate}[label=\arabic*.]
\item Jordan 블록 관점에서 $m_T(t)$에 중근이 생기는 정확한 조건을 서술하십시오.
\item $J_2(2)$에 대해 $J_2(2)^n$을 닫힌식으로 쓰고, $(1,2)$ 성분이 어떻게 증가하는지 설명하십시오.
\item $A=PJP^{-1}$일 때 $e^{tA}=Pe^{tJ}P^{-1}$가 성립하는 이유를 멱급수 정의를 이용해 설명하십시오.
\end{enumerate}

\chapterapplicationtemplate{
상수계수 선형미분방정식
$$
x'(t)=Ax(t),\qquad x(0)=x_0
$$
을 Jordan 표준형으로 풉니다. $A=PJP^{-1}$라 두고 $y(t)=P^{-1}x(t)$로 치환하면
$$
y'(t)=Jy(t),\qquad y(0)=P^{-1}x_0
$$
가 됩니다. $J$가 Jordan 블록들의 직합이므로 각 블록에서 해는
$$
e^{tJ_k(\lambda)}=e^{\lambda t}\sum_{r=0}^{k-1}\frac{t^r}{r!}N^r
$$
형태를 갖습니다. 따라서 전체 해는
$$
x(t)=Pe^{tJ}P^{-1}x_0
$$
로 정리되며, 고유값이 같은 경우에도 다항식 인자 $t^r$가 함께 나타난다는 점이 Jordan 이론의 핵심 응용입니다.
}

\chaptersummarytemplate

\section*{종합 연습문제}

\subsection*{연습문제 A (기초 확인)}
\begin{enumerate}[label=\textbf{A\arabic*.}]
\item $J_4(2)$의 특성다항식과 최소다항식을 구하시오.
\item $J_5(0)$에 대해 $\dim\ker J_5(0)^m$을 $m=1,2,3,4,5$에 대해 계산하시오.
\item 길이 $3$ Jordan 사슬 $(v_1,v_2,v_3)$에서 $T$의 제한사상 행렬을 기저 $(v_1,v_2,v_3)$로 쓰시오.
\item Jordan 블록이 $J_2(1)\oplus J_1(1)\oplus J_3(-1)$일 때, 각 고유값에 대한 기하적 중복도(고유공간 차원)를 구하시오.
\item $J_3(-1)^5$를 이항정리를 이용해 계산하시오.
\item $e^{tJ_2(0)}$를 직접 계산하시오.
\item $\dim\ker N=2,\ \dim\ker N^2=3,\ \dim\ker N^3=5$인 nilpotent $N$의 Jordan 블록 크기를 구하시오.
\item $m_T(t)=(t-3)^2(t+1)$일 때 가능한 Jordan 블록 크기 조합을 모두 쓰시오.
\end{enumerate}

\subsection*{연습문제 B (개념 연결)}
\begin{enumerate}[label=\textbf{B\arabic*.}]
\item nilpotent $N$에 대해 수열 $\dim\ker N,\dim\ker N^2,\dots$가 증가하다가 안정화됨을 증명하시오.
\item $\chi_A(t)=(t-2)^5(t+1)^2$이고
$$
\dim\ker(A-2I)=2,\ \dim\ker(A-2I)^2=4,\ \dim\ker(A-2I)^3=5,
$$
$$
\dim\ker(A+I)=1,\ \dim\ker(A+I)^2=2
$$
일 때 $A$의 Jordan 형식을 구하시오.
\item $A$가 대각화 가능할 필요충분조건이 ``모든 Jordan 블록 크기 $1$''임을 Jordan 형식의 유일성과 연결해 설명하시오.
\item $A=PJP^{-1}$, $p(t)\in F[t]$일 때 $p(A)=Pp(J)P^{-1}$를 증명하고, 이를 이용해 $A^n$ 계산 절차를 정리하시오.
\item $J_k(\lambda)$에 대해 $(J_k(\lambda)-\lambda I)^r$의 rank 공식을 구하시오.
\end{enumerate}

\subsection*{연습문제 C (증명/종합)}
\begin{enumerate}[label=\textbf{C\arabic*.}]
\item 유사한 두 행렬 $A,B$에 대해 $\dim\ker(A-\lambda I)^j=\dim\ker(B-\lambda I)^j$가 모든 $\lambda,j$에 대해 성립함을 증명하시오.
\item $N$이 nilpotent일 때 블록 크기 분할이 $\{\dim\ker N^j\}_{j\ge1}$ 데이터로 유일하게 복원됨을 독립적으로 다시 증명하시오.
\item $A=J_3(2)\oplus J_2(2)$에 대한 미분방정식 $x'(t)=Ax(t)$의 일반해를 구하고, 해에 나타나는 다항식 인자의 차수를 해석하시오.
\end{enumerate}